% !TeX program = pdflatex
\documentclass[11pt,a4paper]{amsart}

\newcommand{\notelanguage}{finnish}
% Shared preamble for course notes and exercise sets.

\usepackage[T1]{fontenc}
\usepackage{lmodern}
\providecommand{\notelanguage}{english}
\usepackage[finnish,english]{babel}
\AtBeginDocument{\expandafter\selectlanguage\expandafter{\notelanguage}}

\newcommand{\notesfinnishlanguage}{finnish}
\ifx\notelanguage\notesfinnishlanguage
\newcommand{\notestheoremname}{Lause}
\newcommand{\noteslemmaname}{Lemma}
\newcommand{\notespropositionname}{Propositio}
\newcommand{\notescorollaryname}{Korollaari}
\newcommand{\notesdefinitionname}{Määritelmä}
\newcommand{\notesexamplename}{Esimerkki}
\newcommand{\notesproblemname}{Tehtävä}
\newcommand{\notesexercisename}{Tehtävä}
\newcommand{\notesremarkname}{Huomautus}
\newcommand{\notessolutionname}{Ratkaisu}
\else
\newcommand{\notestheoremname}{Theorem}
\newcommand{\noteslemmaname}{Lemma}
\newcommand{\notespropositionname}{Proposition}
\newcommand{\notescorollaryname}{Corollary}
\newcommand{\notesdefinitionname}{Definition}
\newcommand{\notesexamplename}{Example}
\newcommand{\notesproblemname}{Problem}
\newcommand{\notesexercisename}{Exercise}
\newcommand{\notesremarkname}{Remark}
\newcommand{\notessolutionname}{Solution}
\fi

\usepackage[a4paper,margin=30mm]{geometry}
\usepackage{microtype}
\usepackage{graphicx}
\usepackage{enumitem}

\newcommand{\exerciseheading}[2]{%
  \par\noindent
  {\large\bfseries #1\par}
  \vspace{0.1em}
  \noindent{\large #2\par}
  \vspace{1.5em}
}

\usepackage{amsmath}
\usepackage{amssymb}
\usepackage{amsthm}
\usepackage{mathtools}
\usepackage{aliascnt}

% Row operations beside the affected row, e.g. \rowop{+2R_1}.
\newcommand{\rowop}[1]{\mathrel{\big|}\,#1}

% Linear systems: the argument gives the number of alignment pairs.
\newenvironment{linearsystem}[1]
{\left\{
  \begin{alignedat}{#1}}
    {
  \end{alignedat}\right.}

  \usepackage[hidelinks,pdfusetitle]{hyperref}

  \numberwithin{equation}{section}
  \setlist{itemsep=0.25em,topsep=0.5em}

  \theoremstyle{plain}
  \newtheorem{theorem}{\notestheoremname}[section]
  \newaliascnt{lemma}{theorem}
  \newtheorem{lemma}[lemma]{\noteslemmaname}
  \aliascntresetthe{lemma}
  \newaliascnt{proposition}{theorem}
  \newtheorem{proposition}[proposition]{\notespropositionname}
  \aliascntresetthe{proposition}
  \newaliascnt{corollary}{theorem}
  \newtheorem{corollary}[corollary]{\notescorollaryname}
  \aliascntresetthe{corollary}

  \theoremstyle{definition}
  \newaliascnt{definition}{theorem}
  \newtheorem{definition}[definition]{\notesdefinitionname}
  \aliascntresetthe{definition}
  \newaliascnt{example}{theorem}
  \newtheorem{example}[example]{\notesexamplename}
  \aliascntresetthe{example}
  \newaliascnt{problem}{theorem}
  \newtheorem{problem}[problem]{\notesproblemname}
  \aliascntresetthe{problem}
  \newtheorem{exercise}{\notesexercisename}

  \theoremstyle{remark}
  \newaliascnt{remark}{theorem}
  \newtheorem{remark}[remark]{\notesremarkname}
  \aliascntresetthe{remark}

  \newenvironment{solution}
  {
  \begin{proof}[\notessolutionname]}
    {
  \end{proof}}

  \usepackage[nameinlink,noabbrev]{cleveref}

  \crefname{theorem}{theorem}{theorems}
  \crefname{lemma}{lemma}{lemmas}
  \crefname{proposition}{proposition}{propositions}
  \crefname{corollary}{corollary}{corollaries}
  \crefname{definition}{definition}{definitions}
  \crefname{example}{example}{examples}
  \crefname{problem}{problem}{problems}
  \crefname{exercise}{exercise}{exercises}
  \crefname{remark}{remark}{remarks}

  \ifx\notelanguage\notesfinnishlanguage
  \crefname{theorem}{lause}{lauseet}
  \crefname{lemma}{lemma}{lemmat}
  \crefname{proposition}{propositio}{propositiot}
  \crefname{corollary}{korollaari}{korollaarit}
  \crefname{definition}{määritelmä}{määritelmät}
  \crefname{example}{esimerkki}{esimerkit}
  \crefname{problem}{tehtävä}{tehtävät}
  \crefname{exercise}{tehtävä}{tehtävät}
  \crefname{remark}{huomautus}{huomautukset}
  \fi

  \newcommand{\NN}{\mathbb{N}}
  \newcommand{\NNone}{\mathbb{N}_1}
  \newcommand{\ZZ}{\mathbb{Z}}
  \newcommand{\QQ}{\mathbb{Q}}
  \newcommand{\RR}{\mathbb{R}}
  \newcommand{\CC}{\mathbb{C}}
  \newcommand{\PP}{\mathbb{P}}

  \DeclareMathOperator{\im}{im}
  \DeclareMathOperator{\rank}{rank}
  \DeclareMathOperator{\syt}{syt}
  \DeclareMathOperator{\pyj}{pyj}
  \DeclareMathOperator{\dom}{dom}
  \DeclarePairedDelimiter{\abs}{\lvert}{\rvert}
  \DeclarePairedDelimiter{\norm}{\lVert}{\rVert}
  \DeclarePairedDelimiterXPP{\maxnorm}[1]{}{\lVert}{\rVert}{_\infty}{#1}
  \newcommand{\set}[1]{\left\lbrace#1\right\rbrace}
  \newcommand{\maxset}[1]{\max\set{#1}}
  \newcommand{\minset}[1]{\min\set{#1}}
  \newcommand{\powerset}[1]{\mathcal{P}\!\left(#1\right)}

\usepackage{tikz}
\usepackage{needspace}

\title{Analyysi II\\Harjoitus 4}
\author{}
\date{}

\begin{document}

\exerciseheading{Analyysi II}{Harjoitus 4}

\begin{exercise}
  Osoita määritelmän perusteella, että funktion
  \(f \colon \RR \to \RR\), \(f(x) = x^2\), derivaatta pisteessä \(2\)
  on \(4\).
\end{exercise}

\begin{solution}
  Suoraan derivaatan määritelmästä saadaan
  \[
    \begin{aligned}
      f^{\prime}(2) &= \lim_{x \to 2} \frac{f(x) - f(2)}{x - 2} \\
      &= \lim_{x \to 2} \frac{x^2 - 4}{x - 2} \\
      &= \lim_{x \to 2} \frac{(x - 2)(x + 2)}{x - 2} \\
      &= \lim_{x \to 2} (x + 2) \\
      &= 4.
    \end{aligned}
  \]
  Siis \(f\) on derivoituva pisteessä \(2\) ja \(f^{\prime}(2) = 4\).
\end{solution}

\begin{exercise}
  Olkoon \(A \subset \RR\) avoin väli ja olkoot
  \(f, g \colon A \to \RR\) funktioita, jotka ovat derivoituvia pisteessä
  \(x_0 \in A\). Osoita, että funktio \(f + g \colon A \to \RR\) on
  derivoituva pisteessä \(x_0\) ja
  \[
    (f + g)'(x_0) = f'(x_0) + g'(x_0).
  \]
\end{exercise}

\begin{solution}
  Kun \(x \in A\) ja \(x \neq x_0\), saadaan
  \[
    \frac{(f + g)(x) - (f + g)(x_0)}{x - x_0}
    = \frac{f(x) - f(x_0)}{x - x_0}
    + \frac{g(x) - g(x_0)}{x - x_0}.
  \]
  Koska \(f\) ja \(g\) ovat derivoituvia pisteessä \(x_0\), oikean puolen
  erotusosamäärillä on äärelliset raja-arvot \(f'(x_0)\) ja \(g'(x_0)\),
  kun \(x \to x_0\). Raja-arvojen summasäännön nojalla siis
  \[
    \begin{aligned}
      \lim_{x \to x_0} \frac{(f + g)(x) - (f + g)(x_0)}{x - x_0}
      &= \lim_{x \to x_0} \frac{f(x) - f(x_0)}{x - x_0}
      + \lim_{x \to x_0} \frac{g(x) - g(x_0)}{x - x_0} \\
      &= f'(x_0) + g'(x_0).
    \end{aligned}
  \]
  Siis \(f + g\) on derivaatan määritelmän nojalla derivoituva pisteessä
  \(x_0\) ja \((f + g)'(x_0) = f'(x_0) + g'(x_0)\).
\end{solution}

\begin{exercise}
  \label{ex:sqrt-abs-origin}
  Olkoon \(f \colon \RR \to \RR\),
  \(f(x) = \sqrt{\abs{x}}\). Tutki määritelmän nojalla, onko funktio
  \(f\) derivoituva origossa.
\end{exercise}

\begin{solution}
  Kun \(x \neq 0\), funktion \(f\) erotusosamäärä origossa on
  \[
    \frac{f(x) - f(0)}{x} = \frac{\sqrt{\abs{x}}}{x}.
  \]
  Olkoon \(M \in \RR\) mielivaltainen ja valitaan
  \(\delta = 1/(M^2 + 1) > 0\). Kun \(0 < x < \delta\), saadaan
  \[
    \frac{f(x) - f(0)}{x}
    = \frac{1}{\sqrt{x}}
    > \frac{1}{\sqrt{\delta}}
    = \sqrt{M^2 + 1} > \abs{M} \geq M.
  \]
  Täten
  \[
    \lim_{x \to 0^+} \frac{f(x) - f(0)}{x} = +\infty.
  \]
  Samalla valinnalla \(\delta = 1/(M^2 + 1)\) saadaan, kun
  \(-\delta < x < 0\),
  \[
    \frac{f(x) - f(0)}{x}
    = -\frac{1}{\sqrt{-x}}
    < -\frac{1}{\sqrt{\delta}}
    = -\sqrt{M^2 + 1} < -\abs{M} \leq M.
  \]
  Siis
  \[
    \lim_{x \to 0^-} \frac{f(x) - f(0)}{x} = -\infty.
  \]
  Erotusosamäärällä ei siis ole äärellistä raja-arvoa origossa, joten
  \(f\) ei ole derivaatan määritelmän nojalla derivoituva origossa.
\end{solution}

\begin{exercise}
  Olkoon \(f \colon \RR \to \RR\) derivoituva.
  Määritellään funktio \(g \colon \RR \to \RR\) asettamalla
  \(g(x) = \abs{f(x)}\). Osoita, että funktio \(g\) on derivoituva
  kaikissa niissä pisteissä \(x \in \RR\), joille \(f(x) \neq 0\).
  Määrää funktion \(g\) derivaatan arvo näissä pisteissä.
\end{exercise}

\begin{solution}
  Olkoon \(x \in \RR\) siten, että \(f(x) \neq 0\).
  Derivoituvana \(f\) on jatkuva pisteessä \(x\). Koska
  \(\abs{f(x)} > 0\), jatkuvuuden nojalla on olemassa \(\delta > 0\)
  siten, että
  \[
    \abs{f(y) - f(x)} < \abs{f(x)},
  \]
  kun \(\abs{y - x} < \delta\). Tällöin
  \[
    f(x) - \abs{f(x)} < f(y) < f(x) + \abs{f(x)}.
  \]
  Jos \(f(x) > 0\), alaraja on \(0\), joten \(f(y) > 0\). Jos taas
  \(f(x) < 0\), yläraja on \(0\), joten \(f(y) < 0\). Siis \(f(y)\)
  on samanmerkkinen kuin \(f(x)\) aina, kun \(\abs{y - x} < \delta\).

  Olkoon \(\varepsilon > 0\). Funktion \(f\) derivoituvuuden nojalla on
  olemassa \(\eta > 0\) siten, että
  \[
    \left|\frac{f(y) - f(x)}{y - x} - f'(x)\right| < \varepsilon,
    \qquad 0 < \abs{y - x} < \eta.
  \]
  Valitaan \(\rho = \min\{\delta, \eta\} > 0\). Kun
  \(0 < \abs{y - x} < \rho\), sekä merkin säilyminen että edellinen
  arvio ovat voimassa.

  Oletetaan ensin, että \(f(x) > 0\). Tällöin \(g(x) = f(x)\) ja
  \(g(y) = f(y)\), kun \(\abs{y - x} < \delta\). Näin ollen
  \[
    \frac{g(y) - g(x)}{y - x}
    = \frac{f(y) - f(x)}{y - x},
    \qquad 0 < \abs{y - x} < \delta.
  \]
  Kun \(0 < \abs{y - x} < \rho\), saadaan siis
  \[
    \left|\frac{g(y) - g(x)}{y - x} - f'(x)\right|
    = \left|\frac{f(y) - f(x)}{y - x} - f'(x)\right|
    < \varepsilon.
  \]
  Koska \(\varepsilon > 0\) oli mielivaltainen, tämä osoittaa, että
  \[
    \lim_{y \to x} \frac{g(y) - g(x)}{y - x}
    = \lim_{y \to x} \frac{f(y) - f(x)}{y - x}
    = f'(x).
  \]
  Täten \(g\) on derivoituva pisteessä \(x\) ja \(g'(x) = f'(x)\).

  Oletetaan sitten, että \(f(x) < 0\). Tällöin \(g(x) = -f(x)\) ja
  \(g(y) = -f(y)\), kun \(\abs{y - x} < \delta\). Näin ollen
  \[
    \frac{g(y) - g(x)}{y - x} =
    -\frac{f(y) - f(x)}{y - x},
    \qquad 0 < \abs{y - x} < \delta.
  \]
  Kun \(0 < \abs{y - x} < \rho\), saadaan vastaavasti
  \[
    \left|\frac{g(y) - g(x)}{y - x} - (-f'(x))\right|
    = \left|\frac{f(y) - f(x)}{y - x} - f'(x)\right|
    < \varepsilon.
  \]
  Koska \(\varepsilon > 0\) oli mielivaltainen, saadaan
  \[
    \lim_{y \to x} \frac{g(y) - g(x)}{y - x}
    = -\lim_{y \to x} \frac{f(y) - f(x)}{y - x}
    = -f'(x).
  \]
  Täten \(g\) on derivoituva pisteessä \(x\) ja \(g'(x) = -f'(x)\).

  Siis \(g\) on derivoituva kaikissa pisteissä \(x\), joissa
  \(f(x) \neq 0\), ja näissä pisteissä
  \[
    g'(x) =
    \begin{cases}
      f'(x), & \text{kun } f(x) > 0, \\
      -f'(x), & \text{kun } f(x) < 0.
    \end{cases}
  \]
\end{solution}

\begin{lemma}[Neliöjuuren jatkuvuus]
  \label{lem:sqrt-continuity}
  Funktio \(x \mapsto \sqrt{x}\) on jatkuva joukossa \([0, \infty)\).
\end{lemma}

\begin{proof}
  Kun \(x, a \geq 0\), saadaan
  \[
    \abs*{\sqrt{x} - \sqrt{a}}^2
    \leq \abs*{\sqrt{x} - \sqrt{a}}(\sqrt{x} + \sqrt{a})
    = \abs{x - a}.
  \]
  Olkoon \(a \geq 0\) ja \(\varepsilon > 0\). Valitaan
  \(\delta = \varepsilon^2 > 0\). Jos \(x \geq 0\) ja
  \(\abs{x - a} < \delta\), niin
  \[
    \abs{\sqrt{x} - \sqrt{a}}
    \leq \sqrt{\abs{x - a}} < \sqrt{\delta} = \varepsilon.
  \]
  Siis neliöjuurifunktio on jatkuva jokaisessa pisteessä \(a \geq 0\).
\end{proof}

\begin{exercise}
  Olkoon \(f \colon [0, \infty) \to \RR\),
  \(f(x) = \sqrt{x}\). Laske määritelmän perusteella funktion \(f\)
  derivaatta \(f'\) pisteessä \(x_0 \in (0, \infty)\). Laske
  \(\lim_{x_0 \to 0^+} f'(x_0)\). Onko funktio \(f\) oikealta
  derivoituva origossa? Selitä havaitsemasi ilmiö geometrisesti.
\end{exercise}

\begin{solution}
  Olkoon \(x_0 \in (0, \infty)\). Lemman~\ref{lem:sqrt-continuity}
  nojalla \(\sqrt{x} \to \sqrt{x_0}\), kun \(x \to x_0\).
  Derivaatan määritelmän ja raja-arvojen laskusääntöjen nojalla saadaan
  \[
    \begin{aligned}
      \lim_{x \to x_0} \frac{f(x) - f(x_0)}{x - x_0}
      &= \lim_{x \to x_0} \frac{\sqrt{x} - \sqrt{x_0}}{x - x_0} \\
      &= \lim_{x \to x_0} \frac{\sqrt{x} - \sqrt{x_0}}{x - x_0}
      \frac{\sqrt{x} + \sqrt{x_0}}{\sqrt{x} + \sqrt{x_0}} \\
      &= \lim_{x \to x_0} \frac{x - x_0}{(x - x_0)(\sqrt{x} + \sqrt{x_0})} \\
      &= \lim_{x \to x_0} \frac{1}{\sqrt{x} + \sqrt{x_0}} \\
      &= \frac{1}{2\sqrt{x_0}}.
    \end{aligned}
  \]
  Siis \(f\) on derivoituva jokaisessa pisteessä \(x_0 > 0\) ja
  \(f'(x_0) = 1/(2\sqrt{x_0})\). Tehtävässä~\ref{ex:sqrt-abs-origin}
  osoitettiin, että \(1/\sqrt{x} \to +\infty\), kun \(x \to 0^+\).
  Siksi myös
  \[
    \lim_{x_0 \to 0^+} f'(x_0)
    = \lim_{x_0 \to 0^+} \frac{1}{2\sqrt{x_0}}
    = +\infty.
  \]
  Saman tehtävässä~\ref{ex:sqrt-abs-origin} todistetun raja-arvon nojalla
  origossa saadaan
  \[
    \lim_{x \to 0^+} \frac{f(x) - f(0)}{x - 0} =
    \lim_{x \to 0^+} \frac{\sqrt{x}}{x} =
    \lim_{x \to 0^+} \frac{1}{\sqrt{x}} = +\infty.
  \]
  Koska erotusosamäärällä ei ole äärellistä raja-arvoa, \(f\) ei ole
  oikealta derivoituva origossa.

  Geometrisesti \(f'(x_0) \to +\infty\) tarkoittaa, että kuvaajan
  tangenttien kulmakertoimet kasvavat rajatta, kun \(x_0 \to 0^+\).
  Myös origon ja kuvaajan pisteen \((x, \sqrt{x})\) kautta kulkevien
  sekanttien kulmakertoimet \(1/\sqrt{x}\) kasvavat rajatta, kun
  \(x \to 0^+\). Sekantit lähestyvät siis pystysuoraa suoraa \(x = 0\),
  joka on kuvaajan tangentti origossa.
\end{solution}

\begin{exercise}
  Olkoot \(r > 0\) ja \(f \colon (-r, r) \to \RR\) sellainen funktio,
  että
  \[
    \abs{f(x)} \leq x^2 \qquad\text{kaikilla } x \in (-r, r).
  \]
  Osoita, että funktio \(f\) on derivoituva origossa ja määrää derivaatta
  \(f'(0)\).
\end{exercise}

\begin{solution}
  Huomataan ensin, että
  \[
    0 \leq \abs{f(0)} \leq 0^2 = 0.
  \]
  Täten \(f(0) = 0\). Olkoon \(\varepsilon > 0\) ja valitaan \(\delta
  = \minset{r, \varepsilon} > 0\). Tällöin
  \[
    \abs*{\frac{f(x) - f(0)}{x - 0}} = \abs*{\frac{f(x)}{x}}
    \leq \frac{x^2}{\abs{x}} = \frac{\abs{x}^2}{\abs{x}} = \abs{x} <
    \varepsilon,
  \]
  kun \(0 < \abs{x - 0} = \abs{x} < \delta\). Täten \(f\) on derivoituva
  origossa ja \(f^{\prime}(0) = 0\).
\end{solution}

\end{document}
